\chapter{Kustomize: Manifest theo bộ}
\label{ch:kustomize}

\section{Vấn đề}

Phần~V sinh ra cả tá file YAML. Apply từng file một thì dễ quên mất
một file; và ngay khi có môi trường thứ hai (``prod, nhưng 2 replica
và secret thật''), copy-paste manifest cho từng môi trường bảo đảm sẽ
trôi (drift) cấu hình. Hai công cụ thống trị mảng này; dự án dùng công
cụ được tích hợp sẵn trong kubectl.

\section{Kustomization, có chú giải}

\code{kubectl apply -k deploy/k8s/base} đọc một file duy nhất:

\begin{lstlisting}[language=yaml]
# deploy/k8s/base/kustomization.yaml
namespace: ts                       (*@\co{1}@*)
resources:                          (*@\co{2}@*)
  - postgres-course.yaml
  - postgres-auth.yaml
  - mongodb.yaml
  - minio.yaml
  # kafka.yaml requires the Strimzi operator (one-time admin step:
  # ../install-strimzi.sh) -- without its CRDs, apply fails on kind Kafka.
  - kafka.yaml
  - kafka-topics.yaml
  - maildev.yaml
  - config-server.yaml
  - discovery-server.yaml
  - course-service.yaml
  - auth-service.yaml
  - dictionary-service.yaml
  - notification-service.yaml
  - api-gateway.yaml
commonLabels:                       (*@\co{3}@*)
  app.kubernetes.io/part-of: teacher-supporter
\end{lstlisting}

\begin{description}
  \item[\co{1}] Đóng dấu lên mọi resource --- không manifest nào cần
  dòng \code{namespace:} riêng, và chuyển cả ứng dụng sang namespace
  khác chỉ là sửa một dòng.
  \item[\co{2}] Danh sách chính thức. Hai sự vắng mặt có chủ ý dạy
  nhiều hơn những gì có mặt: \code{namespace.yaml} bị loại vì nó thuộc
  phạm vi cluster và Role theo namespace của Jenkins không thể apply
  nó (Chương~22); các Secret bị loại vì những lý do của Chương~16. Một
  đối tượng thứ ba cũng ở diện chỉ-admin thì thậm chí không phải là
  file: Strimzi operator và các CRD của nó, được cài một lần bằng
  \code{install-strimzi.sh} (Chương~25) --- \code{kafka.yaml} có trong
  danh sách nhưng không thể apply chừng nào chúng chưa tồn tại. Thứ tự
  liệt kê chỉ mang tính tài liệu --- API server nhận tất cả, còn quá
  trình reconcile lo phần thực tế.
  \item[\co{3}] Một label đóng lên mọi object, biến cả ứng dụng thành
  một selector: \code{kubectl get all -l
  app.kubernetes.io/part-of=teacher-supporter} --- hoặc, trong một
  buổi diễn tập thảm họa, một lệnh delete.
\end{description}

\section{Base và overlay}

Tính năng thật sự của Kustomize xuất hiện cùng môi trường thứ hai.
\emph{Base} giữ sự thật; các \emph{overlay} vá lên nó:

\begin{lstlisting}[language=yaml]
# deploy/k8s/overlays/prod/kustomization.yaml
resources:
  - ../../base
patches:
  - target: { kind: Deployment, name: course-service }
    patch: |-
      - op: replace
        path: /spec/replicas
        value: 2
images:
  - name: localhost:5000/course-service
    newTag: v1.4.2
\end{lstlisting}

Không template, không biến --- overlay là \emph{phép gộp YAML thuần}.
Mỗi lớp đều là YAML Kubernetes hợp lệ mà bạn đọc được, diff được, và
đẩy qua \code{kubectl kustomize} để xem kết quả cuối cùng trước khi
bất cứ thứ gì chạm vào cluster. Bố cục thư mục đã sẵn dạng overlay
(\code{base/}), nên ngày cần overlay prod, đó là một lệnh
\code{mkdir}, không phải một cuộc di dời.

\begin{decision}[Kustomize, không phải Helm]
Helm là câu trả lời còn lại: manifest dưới dạng template Go
(\code{replicas: \{\{ .Values.replicaCount \}\}}) với file giá trị cho
từng môi trường, đóng gói thành các \emph{chart} có phiên bản, chia sẻ
được. Template mạnh hơn hẳn --- vòng lặp, điều kiện, một chart triển
khai được ứng dụng của bất kỳ ai --- và đó chính là điểm rẽ: Helm tỏa
sáng khi bạn \emph{phân phối} phần mềm lên cluster của người lạ; còn
để triển khai ứng dụng của chính mình, template khiến mọi manifest
không đọc nổi cho đến khi được render. Kustomize giữ YAML thuần đúng
là YAML thuần, đi kèm trong kubectl, và phủ hết việc vá
replica/image/label --- toàn bộ nhu cầu ở đây. Bạn vẫn sẽ \emph{dùng}
Helm chart cho phần mềm bên thứ ba (Strimzi, kube-prometheus-stack);
còn viết một chart thì có thể đợi đến khi bạn có người dùng.
\end{decision}

\begin{pitfall}[spec.selector: field is immutable]
Triệu chứng: \code{kubectl apply -k} trong pipeline thất bại ở một
Deployment với lỗi này, dù cùng file đó apply bằng tay thì ổn. Nguyên
nhân: \code{commonLabels} đóng label của nó không chỉ vào metadata mà
vào cả \emph{selector} của mọi Deployment --- và selector là bất biến
sau khi tạo. Một object được tạo lần đầu \emph{bên ngoài} Kustomize
(\code{kubectl apply -f maildev.yaml} lúc thiết lập cột mốc~6,
selector \code{\{app: maildev\}}) sẽ không bao giờ được Kustomize cập
nhật về sau (selector
\code{\{app: maildev, part-of: teacher-supporter\}}). Cách sửa: xóa
object tạo bằng tay và để pipeline tạo lại. Quy tắc: object do một
kustomization quản lý thì \emph{chỉ} apply thông qua nó --- muốn thử
một lần, hãy render bằng \code{kubectl kustomize} rồi pipe kết quả
đó, không bao giờ dùng file thô. (\code{labels:} với
\code{includeSelectors: false} của Kustomize hiện đại tránh được việc
chèn vào selector --- nhưng đổi sang lúc này sẽ phá mọi selector đã
tạo theo cách cũ.)
\end{pitfall}

\begin{pitfall}[hôm qua còn chạy: ảo giác prune]
Triệu chứng: bạn xóa một manifest khỏi \code{resources:}, apply lại,
và object vẫn đang chạy --- hàng tháng trời, cho đến khi nó xung đột
với thứ gì đó. Nguyên nhân: apply-cả-bộ vẫn là apply --- quy tắc của
Chương~14 rằng apply không bao giờ xóa vẫn đúng với Kustomize. Ingress
bị gỡ ở Chương~4 đã cần một lệnh \code{kubectl delete} thủ công. Các
cách sửa có hệ thống: \code{kubectl apply -k --prune} (vẫn bị chặn sau
cờ, và có lý do chính đáng --- một selector gắn nhầm label có thể xóa
hàng loạt), hoặc các controller GitOps (Chương~27) mà toàn bộ công
việc là làm git và cluster hội tụ, kể cả phần xóa.
\end{pitfall}

\begin{handson}[render trước khi tin]
\begin{lstlisting}[language=bash]
kubectl kustomize deploy/k8s/base | less     # the final YAML, merged
kubectl kustomize deploy/k8s/base | grep -c 'kind:'   # object count
kubectl apply -k deploy/k8s/base --dry-run=server     # validate, no changes
\end{lstlisting}
\code{kubectl kustomize} thuần túy là biến đổi văn bản --- đọc đầu ra
của nó là cách bạn xem xét một overlay thật sự thay đổi gì, và
\code{--dry-run=server} là thứ gần nhất với một unit test cho
manifest.
\end{handson}

\begin{quiz}
\begin{enumerate}
  \item Một kustomization giải quyết hai vấn đề gì so với
  \code{kubectl apply -f} từng file, ngay cả khi không có overlay nào?
  \item Vì sao namespace và các Secret cố ý vắng mặt trong
  \code{resources:}? (Hai lý do khác nhau.)
  \item Helm so với Kustomize: nêu sân nhà của mỗi công cụ trong một
  câu, và kể tên phần mềm bên thứ ba trong tương lai của dự án sẽ đến
  dưới dạng chart.
  \item Một đồng nghiệp gỡ \code{mongodb.yaml} khỏi danh sách để
  ``cho nghỉ hưu'' nó. Điều gì thật sự xảy ra, và hai quy trình đúng
  là gì?
\end{enumerate}
\end{quiz}
